\documentclass[12pt]{report}

\usepackage[spanish]{babel}
\usepackage[utf8]{inputenc}
\usepackage[T1]{fontenc}

\usepackage{fancyhdr}
\usepackage{graphicx}
\usepackage{geometry}
\usepackage{titlesec}
\usepackage{xcolor}
\usepackage{cite}
\usepackage{parskip}

\geometry{a4paper, margin=1in}

\definecolor{highlight}{RGB}{255,255,0}
\definecolor{azulito}{RGB}{85,142,213}

\titleformat{\section}
  {\normalfont\Large\bfseries\MakeUppercase}
  {}{0pt}{}

\addto\captionsspanish{\renewcommand{\contentsname}{Contenido}}

\renewcommand{\thesection}{}
\renewcommand{\thesubsection}{}

\begin{document}

% PORTADA

\thispagestyle{empty}

\vspace*{2cm}

\begin{center}

\noindent\rule{\textwidth}{0.4pt}

\vspace{0.5cm}

{\Huge\bfseries ANÁLISIS DE LA VARIACIÓN EN LA CALIDAD DEL AIRE Y SUS EFECTOS SOBRE LA SALUD EN ESTADOS UNIDOS}

\vspace{0.5cm}

\noindent\rule{\textwidth}{0.4pt}

\vspace{0.5cm}

{\large Proyecto de Métodos Numéricos}

\vspace{2cm}

\textbf{ASIGNATURA:} Métodos Numéricos

\vspace{0.5cm}

\textbf{INTEGRANTES:}

GUEVARA JUSTIN, CAHUATE DEREK, LOPEZ ELIZABETH

\vspace{2cm}

\textbf{FECHA DE ENTREGA:}

Junio 2026

\end{center}

\newpage

% ÍNDICE

\tableofcontents
\thispagestyle{empty}

\newpage

% CONTENIDO

\pagenumbering{arabic}
\setcounter{page}{1}

\section{Introducción}

En este proyecto se analiza la variación de la calidad del aire en Estados Unidos usando datos históricos del índice AQI. El objetivo principal es observar cómo cambia este índice a lo largo de los años y comparar su comportamiento por regiones.

\section{Objetivos}

\subsection{Objetivo general}

Analizar la variación de la calidad del aire en Estados Unidos mediante datos históricos, aplicando limpieza de datos, agrupación regional, normalización porcentual e interpolación spline.

\subsection{Objetivos específicos}

\begin{itemize}
    \item Limpiar el conjunto de datos eliminando registros incompletos.
    \item Agrupar los estados por regiones del Censo de Estados Unidos.
    \item Calcular porcentajes para comparar mejor los tipos de días y contaminantes.
    \item Aplicar interpolación spline para observar la tendencia del AQI.
\end{itemize}

\section{Metodología}

Primero se trabajó con el archivo original de datos. Luego se eliminaron filas con valores nulos para evitar errores en los cálculos. Después, cada estado fue asignado a una región: Northeast, Midwest, South o West.

Posteriormente, los datos se agruparon por región y año usando el promedio de las variables seleccionadas. También se agregaron columnas porcentuales para comparar los valores de forma proporcional.

\section{Normalización porcentual}

La normalización porcentual permite comparar datos aunque no todos los años o regiones tengan la misma cantidad de registros.

La fórmula usada fue:

\[
Porcentaje = \frac{Valor}{Total} \times 100
\]

Por ejemplo, para los días buenos:

\[
Good\ Days\_pct = \frac{Good\ Days}{Total\_Dias} \times 100
\]

Esto no elimina los datos originales. Solo agrega nuevas columnas con el sufijo \texttt{\_pct}.

\section{Interpolación spline}

La interpolación spline es un método numérico que permite crear una curva suave a partir de varios puntos conocidos. En este proyecto se usa para observar la tendencia del AQI a través de los años.

\section{Conclusiones}

La limpieza y agrupación de datos permitió ordenar la información por región y año. La normalización porcentual ayudó a comparar los datos de una forma más justa, ya que no todos los registros tienen la misma cantidad de días. Finalmente, el método spline permite visualizar mejor la tendencia del AQI.

% REFERENCIAS

\renewcommand{\bibname}{\MakeUppercase{REFERENCIAS}}

\bibliographystyle{IEEEtran}
\bibliography{referencias}

\end{document}